\ifdefined\uploadfont
    % 使用上传字体方案（需要在 font/ 目录下放好对应文件）
    \usepackage{xeCJK}
    \usepackage{zhnumber}

    % =====================================================
    % 切换字体：取消注释其中一个方案，其余保持注释
    % =====================================================

    % 方案B: 霞鹜新致宋（开源，古典明朝风格）
    \def\cjkMainFont{LXGWNeoZhiSong.ttf}
    \def\cjkBoldFont{SourceHanSerifCN-SemiBold.otf}

    % % 方案C: 宋体-简（macOS 自带，多字重完整）
    % \def\cjkSysFont{}           % 系统字体标记，告知下方用名称而非路径加载
    % \def\cjkMainFont{Songti SC Light}
    % \def\cjkBoldFont{Songti SC Bold}

    % =====================================================
    \def\cjkMonoFont{FandolKai-Regular.otf}
    \def\latinMainFont{EBGaramond.ttf}
    \def\latinItalicFont{EBGaramond.ttf}

    % 覆盖 xeCJK 默认 Script=CJK（Fandol 字体用 script=hani）
    \defaultCJKfontfeatures{RawFeature={script=hani}}

    % 缺字自动回退（用于 \kaiti 等字体缺少罕见字形时兜底）
    \xeCJKsetup{AutoFallBack = true}

    % 主中文字体（系统字体按名称加载；文件字体按路径加载）
    \ifdefined\cjkSysFont
        \setCJKmainfont[
            BoldFont = \cjkBoldFont,
            BoldItalicFont = \cjkBoldFont
        ]{\cjkMainFont}
    \else
        \setCJKmainfont[
            Path = font/,
            UprightFont = *,
            ItalicFont = *,
            BoldFont = \cjkBoldFont,
            BoldItalicFont = \cjkBoldFont
        ]{\cjkMainFont}
    \fi

    % 英文字体
    \setmainfont[
        Path = font/,
        UprightFont = *,
        ItalicFont = *,
        BoldFont = *,
        BoldItalicFont = *,
        AutoFakeSlant = 0.2,
        BoldFeatures     = {RawFeature = {+wght=700}},
        BoldItalicFeatures = {RawFeature = {+wght=700}}
    ]{\latinMainFont}

    % 楷体（FandolKai，GPLv3 + Font Exception；只有 Regular，伪造粗/斜）
    \newCJKfontfamily\kaiti[
        Path = font/,
        RawFeature = {script=hani},
        AutoFakeBold  = 2,
        ItalicFont = *,
        BoldFont = *,
        BoldItalicFont = *
    ]{\cjkMonoFont}

    % 等宽字体
    \setCJKmonofont[
        Path = font/,
        AutoFakeBold  = 2,
        UprightFont = *,
        BoldFont = *,
        ItalicFont = *,
        BoldItalicFont = *
    ]{\cjkMonoFont}

\else
    % 默认使用 ctex 提供的中文字体配置
    \usepackage{ctex}
\fi

% 引用环境字体
\newcommand{\quotefont}{\kaiti}
